\section{优化策略与复杂度分析}
\subsection{策略设计}
优化流程如图~\ref{fig:loop} 所示。其核心不是仅提高单个专家的并发数，而是在轨迹统计下权衡互斥性、热点副本、物理碎片与调度等待。共现低的专家可共享资源候选，频率高且压力大的专家得到更多副本，热点组优先分配到不同 Sub-Cube，并将常用权重偏向低 z 层。该设计的收益应由具体轨迹验证；它不是对任意模型与任意硬件的普适承诺。

\begin{figure}[t]
\centering
\begin{tikzpicture}[
  node distance=4mm,
  box/.style={draw,rounded corners=2pt,minimum width=31mm,minimum height=7mm,align=center,fill=green!7},
  arrow/.style={-{Stealth[length=1.6mm]},thick}
]
\node[box] (a) {轨迹统计：共现、频率、转移};
\node[box,below=of a] (b) {互斥分组与热点选择};
\node[box,below=of b] (c) {精确分配、局部搜索与退火};
\node[box,below=of c] (d) {副本放置与容量检查};
\node[box,below=of d] (e) {静态调度与约束验证};
\draw[arrow] (a)--(b);
\draw[arrow] (b)--(c);
\draw[arrow] (c)--(d);
\draw[arrow] (d)--(e);
\draw[arrow] (e.east) to[out=0,in=0,looseness=1.2] node[right,font=\scriptsize,align=left]{指标反馈\\与参数调优} (a.east);
\end{tikzpicture}
\caption{轨迹感知映射与验证闭环。}
\label{fig:loop}
\end{figure}

\subsection{复杂度}
令 $C$ 为 Weight-Cube 数，$R$ 为 Shelf 行数，$E$ 为专家数，$G$ 为分组数，$I$ 为局部搜索迭代数，$S$ 为退火步数。ShelfPacker 的时间复杂度为 $O(CR)$；多起点互斥分组约为 $O(TE^2G)$，其中 $T$ 是分组试验数；热点子图精确分配最坏为 $O(N^k)$，实际由 branch-and-bound 剪枝；局部搜索为 $O(I(GN+G^2))$；退火为 $O(SG^2)$。调度阶段对每个任务比较候选 placement 并使用优先队列维护就绪集，复杂度约为 $O(Q(P+\log Q))$，其中 $Q$ 为任务数、$P$ 为候选物理副本数。

为避免搜索成本无界，项目将精确求解限制在热点 top-$k$，并提供并行多起点与确定性开关。实际提交参数将 $k$ 固定为 6，使 $N=3$ 时的名义枚举上界为 $9^6$；这使优化重点集中在最可能影响延迟的热点组。
